\chapter{Dandruff}

\section*{Definition and Overview}
Dandruff is a common, non-contagious scalp condition characterised by the shedding of white or grey flakes of dead skin from the scalp, often accompanied by mild itching. It is a mild form of seborrhoeic dermatitis, a chronic inflammatory condition affecting areas of skin rich in sebaceous glands. Dandruff is a cosmetic nuisance for most people rather than a serious medical problem, and it responds well to regular use of medicated shampoos. This chapter covers the causes, features, diagnosis, and treatment of dandruff.

\section*{Epidemiology}
Dandruff is one of the most common scalp conditions worldwide, affecting up to half of adults at some point in their lives. It typically begins after puberty, when sebum production increases, and is most prevalent in young and middle-aged adults. It affects males more often than females and tends to fluctuate with seasons, stress, and hair-care habits. Seborrhoeic dermatitis in its more extensive form affects around 3--5\% of the general population and is more common in people with Parkinson's disease, HIV infection, or other conditions affecting the immune system.

\section*{Aetiology and Pathophysiology}
\begin{itemize}
    \item \textbf{Malassezia yeasts:} commensal fungi that live on the scalp metabolise sebum; in susceptible people an overgrowth triggers an inflammatory response and accelerated skin cell turnover
    \item \textbf{Sebum production:} the oily secretion of sebaceous glands provides the substrate for Malassezia and is central to the condition
    \item \textbf{Individual susceptibility:} genetic factors and immune reactivity determine who develops dandruff and how severe it is
    \item \textbf{Aggravating factors:} stress, cold dry weather, infrequent washing, harsh hair products, and certain neurological or immune conditions
    \item \textbf{Accelerated shedding:} inflammation causes skin cells to be produced and shed more rapidly, producing visible clumps of flakes
\end{itemize}

\section*{Clinical Features}
\begin{itemize}
    \item White or grey flakes on the scalp, hair, and shoulders
    \item Itchy scalp, which may be mild or troublesome
    \item Dry or greasy appearance of the flakes depending on the individual
    \item Redness and scaling along the hairline, eyebrows, or behind the ears in seborrhoeic dermatitis
    \item Symptoms that fluctuate, often worse in winter or during stress
    \item No scarring or permanent hair loss in ordinary dandruff
\end{itemize}

\section*{Differential Diagnosis}
\begin{itemize}
    \item Seborrhoeic dermatitis, of which dandruff is the mild scalp-limited form
    \item Dry scalp from over-washing or harsh products, which sheds fine dry flakes without greasiness
    \item Psoriasis of the scalp, which produces thicker, silvery plaques
    \item Tinea capitis (ringworm of the scalp), more common in children, with patchy scaling and hair loss
    \item Contact dermatitis from hair dyes or styling products
    \item Eczema (atopic dermatitis) involving the scalp in children
\end{itemize}

\section*{When to Seek Medical Care}
Most dandruff can be managed with over-the-counter medicated shampoos. See a doctor if the scaling is severe, spreads to the face or body, fails to improve with treatment, or is accompanied by hair loss, redness, or open sores, as these may indicate another condition such as psoriasis or a fungal infection.

\textbf{Call 000 (or local emergency services) for:}
\begin{itemize}
    \item Any signs of a severe infection such as spreading redness, pain, fever, or discharge from the scalp
\end{itemize}

\section*{Diagnosis and Investigations}
\begin{itemize}
    \item \textbf{Clinical examination:} dandruff is diagnosed from the appearance and distribution of the scaling
    \item \textbf{History:} duration, response to previous treatments, and aggravating factors
    \item \textbf{Skin scrapings:} may be examined under a microscope if tinea capitis is suspected, particularly in children
    \item \textbf{Referral:} to a dermatologist for persistent or atypical cases
\end{itemize}

\section*{Management}
\begin{itemize}
    \item \textbf{Medicated shampoos:} rotate products containing zinc pyrithione, ketoconazole, selenium sulphide, or salicylic acid
    \item \textbf{Application technique:} massage into the scalp and leave for 3--5 minutes before rinsing to allow the active ingredient to work
    \item \textbf{Frequency:} use two to three times weekly at first, then once weekly or as needed to maintain control
    \item \textbf{Coal tar shampoo:} an option for stubborn cases, though it may discolour light hair
    \item \textbf{Topical corticosteroids:} short courses of scalp lotion or foam for inflamed seborrhoeic dermatitis
    \item \textbf{Lifestyle measures:} regular washing, gentle brushing, and stress management
\end{itemize}

\section*{Complications}
\begin{itemize}
    \item Persistent self-consciousness and reduced quality of life related to visible flakes
    \item Scalp irritation or contact dermatitis from excessive use of strong shampoos
    \item Secondary bacterial infection from scratching
    \item Extension of seborrhoeic dermatitis to the face, chest, or other body folds
\end{itemize}

\section*{Prognosis and Outcomes}
Dandruff is a chronic but highly manageable condition. With regular use of an appropriate medicated shampoo, most people achieve good control of flakes and itching, although symptoms commonly recur if treatment is stopped. Seborrhoeic dermatitis tends to run a relapsing course requiring maintenance therapy. Dandruff does not damage the hair or scalp permanently and resolves or becomes negligible in most people with consistent care.

\section*{Prevention and Self-Care}
\begin{itemize}
    \item Wash hair regularly with a gentle shampoo to remove excess oil and flakes
    \item Use an anti-dandruff shampoo preventively once weekly
    \item Avoid harsh styling products and vigorous scratching
    \item Manage stress, which can trigger flare-ups
    \item Protect the scalp from extremes of cold, dry air
    \item See a doctor if symptoms persist despite several weeks of medicated shampoo use
\end{itemize}

\section*{Sources and Further Reading}
The following reputable sources provide further information for healthcare students and patients seeking reliable, current advice about dandruff and its management.

\begin{itemize}
    \item NHS. \textit{Dandruff: causes, treatment, and prevention.} https://www.nhs.uk/conditions/dandruff/
    \item Mayo Clinic. \textit{Dandruff: symptoms and treatment options.} https://www.mayoclinic.org/diseases-conditions/dandruff/
    \item MSD Manual. \textit{Seborrhoeic dermatitis: professional overview.} https://www.msdmanuals.com/
    \item MedlinePlus. \textit{Dandruff health information.} https://medlineplus.gov/dandruff.html
    \item American Academy of Dermatology. \textit{How to treat dandruff.} https://www.aad.org/
    \item DermNet. \textit{Dandruff and seborrhoeic dermatitis.} https://dermnetnz.org/
\end{itemize}
